\lecture{Лекция 5}{}

\subsection{Уточнённые оценки шага и скорости}

Теорема~\ref{thm:grad-rate} оставила коэффициент сжатия в виде
$q = 1 - \eps\,\alpha_k\, m\big(1 + \frac{m}{M}\big)$ — через сам шаг $\alpha_k$.
Доведём выкладку до оценки, в которой $\alpha_k$ уже не участвует.

\subsubsection*{Одношаговая оценка через гессиан}

Запишем приращение за итерацию по формуле Тейлора с остаточным членом в форме
Лагранжа:
\begin{align*}
  f(x_{k+1}) - f(x_k)
  &= f\big(x_k - \alpha_k \nabla f(x_k)\big) - f(x_k) \\
  &= f(x_k) - \alpha_k \nabla\T f(x_k)\,\nabla f(x_k)
     + \tfrac{1}{2}\alpha_k^2\, \nabla\T f(x_k)\, H(\bar{x}_k)\, \nabla f(x_k)
     - f(x_k) ,
\end{align*}
где
\[
  \bar{x}_k = x_k + \theta\,(x_{k+1} - x_k)
            = x_k - \theta\,\alpha_k \nabla f(x_k) , \qquad \theta \in (0,1) .
\]
Подставляя верхнюю границу гессиана $z\T H z \le M\norm{z}^2$, получаем
\begin{equation}\label{eq:one-step-M}
  f(x_{k+1}) - f(x_k)
  \ \le\ -\alpha_k \norm{\nabla f(x_k)}^2
          + \frac{M}{2}\alpha_k^2 \norm{\nabla f(x_k)}^2
  = \norm{\nabla f(x_k)}^2 \Big( \frac{M}{2}\alpha_k^2 - \alpha_k \Big) .
\end{equation}

\begin{idea}
Это та же лемма~\ref{lem:grad-main}, но в более сильной форме. Там квадратичный
член оценивался через липшицеву константу градиента и давал $R\alpha^2$; здесь
используется разложение второго порядка, и коэффициент получается вдвое меньше:
$\frac{M}{2}\alpha^2$. Для дважды дифференцируемой $f$ константа Липшица
градиента как раз равна наибольшему собственному числу гессиана, $R = M$,
так что~\eqref{eq:one-step-M} строго точнее~\eqref{eq:grad-main}.
Поэтому и пороги для шага дальше получаются вдвое мягче:
$\frac{2}{M}(1-\eps)$ вместо $\frac{1-\eps}{R}$.
\end{idea}

\subsubsection*{Границы для принятого шага}

Напомним условие приёма шага~\eqref{eq:armijo}:
\[
  f(x_{k+1}) - f(x_k) \ \le\ -\eps\, \alpha_k \norm{\nabla f(x_k)}^2 ,
  \qquad \eps \in (0,1) .
\]
Из~\eqref{eq:one-step-M} видно, что для его выполнения достаточно
\[
  -\alpha_k + \frac{M}{2}\alpha_k^2 \ \le\ -\eps\,\alpha_k
  \qquad\Longrightarrow\qquad
\]
\begin{equation}\label{eq:alpha-upper-M}
  \alpha_k \ \le\ \frac{2}{M}\,(1 - \eps) .
\end{equation}

\begin{remark}
Алгоритм дробления: $\alpha_{k+1} = \alpha_k \lambda$, где $\lambda \in (0,1)$ —
параметр дробления. Оценки при этом остаются справедливы.
\end{remark}

\begin{idea}[title={Почему дробление не портит оценки}]
Повтор рассуждения шага~4 теоремы~\ref{thm:grad-split}, но уже с константой $M$.
Шаг, \emph{не} прошедший проверку, обязан нарушать~\eqref{eq:alpha-upper-M}:
\[
  -\eps\,\alpha < -\alpha + \frac{\alpha^2}{2}M
  \qquad\Longleftrightarrow\qquad
  \alpha > \frac{2}{M}(1 - \eps) .
\]
Последним отвергнутым был шаг $\alpha_k / \lambda$, поэтому
\begin{equation}\label{eq:alpha-lower-M}
  \frac{\alpha_k}{\lambda} > \frac{2}{M}(1 - \eps)
  \qquad\Longrightarrow\qquad
  \alpha_k > \frac{2}{M}\,\lambda\,(1 - \eps) .
\end{equation}
Итого принятый шаг зажат между двумя положительными константами:
сверху~\eqref{eq:alpha-upper-M} гарантирует, что дробление конечно,
снизу~\eqref{eq:alpha-lower-M} — что убывание функции за итерацию отделено от
нуля. Это и значит «оценки справедливы».
\end{idea}

\subsubsection*{Коэффициент сжатия без $\alpha_k$}

Подставим нижнюю границу шага и оценку~\eqref{eq:rate-5} в $q$ из
теоремы~\ref{thm:grad-rate}:
\[
  q = 1 - \eps\,\alpha_k\, m\Big(1 + \frac{m}{M}\Big)
    \ \le\ 1 - \eps(1-\eps)\,\frac{2m}{M}\Big(1 + \frac{m}{M}\Big) .
\]
Выражение $\eps(1-\eps)$ максимально при $\eps = \tfrac12$, где равно $\tfrac14$,
поэтому наилучшая оценка —
\begin{equation}\label{eq:q-min}
  \eps_{\max} = \frac{1}{2} ,
  \qquad
  q_{\min} = 1 - \frac{1}{2}\,\frac{m}{M}\Big(1 + \frac{m}{M}\Big) ,
  \qquad 0 < \frac{m}{M} \le 1 .
\end{equation}

\begin{idea}[title={Почему $0 \le q_{\min} < 1$}]
Обозначим $t = m/M$. Из $0 < m \le M$ следует $t \in (0,1]$, и
\[
  q_{\min} = 1 - \tfrac{1}{2}\,t\,(1 + t) .
\]
Функция $t \mapsto \tfrac12 t(1+t)$ строго возрастает на $(0,1]$ от $0$ до $1$,
значит $q_{\min} \in [0, 1)$ — сходимость геометрическая. Крайние случаи:
\begin{itemize}[nosep,left=0pt]
  \item $t = 1$ (то есть $m = M$, линии уровня — сферы): $q_{\min} = 0$ —
        метод попадает в минимум за один шаг;
  \item $t \to 0$ (сильно вытянутые линии уровня): $q_{\min} \to 1$ —
        сходимость сколь угодно медленная.
\end{itemize}
Множитель $\lambda$ из~\eqref{eq:alpha-lower-M} здесь опущен; строго говоря
$q = 1 - \eps(1-\eps)\frac{2m\lambda}{M}\big(1 + \frac{m}{M}\big)$, и
записанная оценка отвечает предельному случаю $\lambda \to 1$ (дробление мелкое,
принятый шаг почти максимально допустимый).
\end{idea}

Дальше всё как в шаге~6 теоремы~\ref{thm:grad-rate}:
\[
  f(x_k) - f(x_*) \ \le\ q\,\big(f(x_{k-1}) - f(x_*)\big)
  \ \le\ \dots \ \le\ q^{k}\big(f(x_0) - f(x_*)\big) ,
\]
а вместе с левой частью~\eqref{eq:rate-2}
\[
  \frac{m}{2}\norm{x_k - x_*}^2 \ \le\ f(x_k) - f(x_*)
  \ \le\ q^{k}\big(f(x_0) - f(x_*)\big)
  \quad\Longrightarrow\quad
  \norm{x_k - x_*} \le C\, q^{k/2} ,
\]
\[
  C = \sqrt{\frac{2}{m}\big(f(x_0) - f(x_*)\big)} .
\]

\subsection{Скорость наискорейшего спуска}

\begin{theorem}\label{thm:steepest-rate}
Для метода наискорейшего спуска справедливы оценки
\[
  f(x_{k+1}) - f(x_*) \ \le\ q\,\big(f(x_k) - f(x_*)\big) ,
  \qquad
  \norm{x_k - x_*} \ \le\ C\, q^{k/2} , \quad C < +\infty ,
\]
\[
  q = 1 - \frac{m}{2M}\Big(1 + \frac{m}{M}\Big) .
\]
\end{theorem}

\begin{proof}
Пусть $\alpha_k$ — шаг, полученный по методу наискорейшего спуска, то есть
$\alpha_k = \argmin_{\alpha \ge 0} f\big(x_k - \alpha \nabla f(x_k)\big)$,
а $\alpha$ — любое. По определению минимума и оценке~\eqref{eq:one-step-M},
записанной для произвольного $\alpha$,
\begin{align*}
  f(x_{k+1}) - f(x_k)
  &= f\big(x_k - \alpha_k \nabla f(x_k)\big) - f(x_k)
   \ \le\ f\big(x_k - \alpha \nabla f(x_k)\big) - f(x_k) \\
  &\le \Big(-\alpha + \frac{M}{2}\alpha^2\Big) \norm{\nabla f(x_k)}^2 .
\end{align*}
Неравенство верно при любом $\alpha$, поэтому возьмём лучшее:
\[
  -\alpha + \frac{M}{2}\alpha^2 \ \to\ \min_{\alpha > 0}
  \qquad\Longrightarrow\qquad
  \alpha_* = \frac{1}{M} , \qquad
  -\alpha_* + \frac{M}{2}\alpha_*^2 = -\frac{1}{2M} .
\]
Отсюда
\[
  f(x_{k+1}) - f(x_k) \ \le\ -\,\frac{\norm{\nabla f(x_k)}^2}{2M} ,
\]
а с учётом~\eqref{eq:rate-5}
\[
  f(x_{k+1}) - f(x_k)
  \ \le\ -\,\frac{1}{2M}\Big(1 + \frac{m}{M}\Big) m \big(f(x_k) - f(x_*)\big) .
\]
Наконец,
\[
  f(x_{k+1}) - f(x_*) = \big(f(x_{k+1}) - f(x_k)\big) + \big(f(x_k) - f(x_*)\big)
  \le \Big(1 - \frac{m}{2M}\Big(1 + \frac{m}{M}\Big)\Big)
      \big(f(x_k) - f(x_*)\big) ,
\]
то есть $q = 1 - \frac{m}{2M}\big(1 + \frac{m}{M}\big)$. Оценка по аргументу
получается из~\eqref{eq:rate-2} так же, как выше:
\[
  \norm{x_k - x_*} \le q^{k/2}\, C_* ,
  \qquad
  C_* = \sqrt{\frac{2}{m}\big(f(x_0) - f(x_*)\big)} .
  \qedhere
\]
\end{proof}

\begin{idea}[title={Почему $q$ совпало с $q_{\min}$ дробления}]
Это не совпадение. В методе дробления лучшее значение достигалось при
$\eps = \tfrac12$, а допустимый шаг при таком $\eps$ —
$\alpha \le \frac{2}{M}\big(1 - \tfrac12\big) = \frac{1}{M}$. Ровно это же
значение $\alpha_* = 1/M$ минимизирует квадратичную мажоранту
$-\alpha + \frac{M}{2}\alpha^2$ в доказательстве выше. Обе выкладки
оптимизируют одну и ту же верхнюю оценку, поэтому и дают один и тот же $q$.

Содержательно: наискорейший спуск по построению не хуже любого фиксированного
шага, в том числе $\alpha = 1/M$; оценка для него поэтому наследуется от оценки
для шага $1/M$. Реальная скорость наискорейшего спуска может быть лучше, но
доказать это через ту же мажоранту нельзя.
\end{idea}

\subsubsection*{Ортогональность соседних градиентов}

\begin{idea}[title={Откуда этот материал}]
На лекции зигзагообразная траектория наискорейшего спуска была нарисована
(рис.~\ref{fig:zigzag}), но не обоснована. Утверждение ниже объясняет её
точно: соседние градиенты перпендикулярны.
\end{idea}

\begin{theorem}\label{thm:steepest-orthogonal}
Пусть в методе наискорейшего спуска шаг $\alpha_k$ найден точно и минимум по
$\alpha$ достигается во внутренней точке ($\alpha_k > 0$). Тогда
\[
  \nabla\T f(x_{k+1})\,\nabla f(x_k) = 0 .
\]
\end{theorem}

\begin{proof}
Обозначим $g_k = \nabla f(x_k)$ и введём функцию одной переменной
\[
  \varphi(\alpha) = f\big(x_k - \alpha g_k\big) .
\]
По определению наискорейшего спуска $\alpha_k = \argmin_{\alpha > 0}
\varphi(\alpha)$, и если минимум внутренний, то $\varphi'(\alpha_k) = 0$. По
правилу дифференцирования сложной функции
\[
  \varphi'(\alpha) = \nabla\T f\big(x_k - \alpha g_k\big)\,(-g_k) ,
\]
а так как $x_{k+1} = x_k - \alpha_k g_k$, подстановка $\alpha = \alpha_k$ даёт
\[
  0 = \varphi'(\alpha_k) = -\nabla\T f(x_{k+1})\, g_k
  \qquad\Longrightarrow\qquad
  \nabla\T f(x_{k+1})\,\nabla f(x_k) = 0 . \qedhere
\]
\end{proof}

\begin{idea}
Отсюда и берётся зигзаг: направление очередного шага $-\nabla f(x_{k+1})$
перпендикулярно предыдущему, то есть траектория всё время поворачивает ровно на
прямой угол. На вытянутых линиях уровня (большое $M/m$) такие повороты очень
мелко «пилят» вдоль оврага, что и объясняет медленную сходимость с
$q \to 1$ из теоремы~\ref{thm:steepest-rate}.

Содержательно условие $\varphi'(\alpha_k) = 0$ означает: двигаться вдоль
текущего направления больше нечего — вся информация о нём исчерпана. Но на
следующем шаге метод про него забывает, и часть продвижения приходится
отыгрывать заново. Именно это исправляют методы сопряжённых градиентов
(\S\ref{subsec:conjugate-gradients}): там соседние направления делают
ортогональными не в обычном смысле, а в метрике гессиана, и для квадратичной
функции зигзаг исчезает вовсе.
\end{idea}

Итак, оценки получены и для шага, и для сходимости, и для её скорости. Отсюда
естественный вопрос: нельзя ли взять \emph{постоянный} шаг и не решать на каждой
итерации одномерную задачу?
